% Chapter 3: Tangent Vectors
% Lee, Introduction to Smooth Manifolds (2nd ed., GTM 218).
% Generated from the section extraction; nodes carry only Lean that is
% verified sorry-free. See ../../UPSTREAM_LEAN_AUDIT.md.


\section{Tangent Vectors}

\begin{definition}
\label{def:3-13-extra-1}
\lean{geometric_tangent_space}
\leanok


For $a\in\mathbb{R}^n$, the geometric tangent space at $a$ is
\[
\mathbb{R}^n_a=\{a\}\times\mathbb{R}^n
  =\{(a,v):v\in\mathbb{R}^n\}.
\]
Write $(a,v)$ as $v_a$, or as $v|_a$ when needed. Its vector-space operations are
\[
v_a+w_a=(v+w)_a,\qquad c(v_a)=(cv)_a,
\]
and $(e_i|_a)_{i=1}^n$ is a basis. A geometric tangent vector is an element of
some $\mathbb{R}^n_a$; the base-point index makes the spaces at distinct points
disjoint.
\end{definition}

\begin{definition}
\label{def:3-13-extra-2}
\lean{PointDerivation}

\mathlibok
For $a\in\mathbb{R}^n$, a derivation at $a$ is an $\mathbb{R}$-linear map
$w:C^\infty(\mathbb{R}^n)\to\mathbb{R}$ satisfying
\[
w(fg)=f(a)w(g)+g(a)w(f). \tag{3.3}
\]
The vector space of such derivations is denoted $T_a\mathbb{R}^n$, with
pointwise operations
\[
(w_1+w_2)(f)=w_1(f)+w_2(f),qquad (cw)(f)=c\,w(f).
\]
\end{definition}

\begin{lemma}[Properties of Derivations]
\label{lem:3-1}
\lean{point_derivation_mul_eq_zero_of_vanish_at}
\leanok


\dcref{ch3:3.1}
Let $a\in\mathbb{R}^n$, $w\in T_a\mathbb{R}^n$, and
$f,g\in C^\infty(\mathbb{R}^n)$. Then:
\begin{enumerate}
\item if $f$ is constant, then $w(f)=0$;
\item if $f(a)=g(a)=0$, then $w(fg)=0$.
\end{enumerate}
\end{lemma}

\begin{proof}
\leanok
\sketch
For the constant function $1$, the product rule gives
$w(1)=w(1\cdot1)=2w(1)$, hence $w(1)=0$; linearity handles every constant.
If $f(a)=g(a)=0$, (3.3) gives $w(fg)=0$ directly.
\end{proof}

\begin{definition}
\label{def:3-13-extra-3}
\lean{PointDerivation}

\mathlibok
Let $M$ be a smooth manifold, possibly with boundary, and let $p\in M$. A
derivation at $p$ is an $\mathbb{R}$-linear map
$v:C^\infty(M)\to\mathbb{R}$ such that
\[
v(fg)=f(p)v(g)+g(p)v(f)qquad
\text{for all }f,g\in C^\infty(M). \tag{3.4}
\]
The vector space of derivations at $p$ is the tangent space $T_pM$; its
elements are tangent vectors at $p$.
\end{definition}

\begin{lemma}[Properties of Tangent Vectors on Manifolds]
\label{lem:3-4}
\lean{tangent_vector_mul_eq_zero_of_vanish_at}
\leanok


\dcref{ch3:3.4}
Let $M$ be a smooth manifold, possibly with boundary, let $p\in M$ and
$v\in T_pM$, and let $f,g\in C^\infty(M)$. Then:
\begin{enumerate}
\item if $f$ is constant, then $v(f)=0$;
\item if $f(p)=g(p)=0$, then $v(fg)=0$.
\end{enumerate}
\end{lemma}



\section{The Differential of a Smooth Map}

\begin{definition}
\label{def:3-14-extra-1}
\lean{mfderiv}

\mathlibok
For a smooth map $F:M\to N$ between smooth manifolds, possibly with boundary,
and $p\in M$, the differential at $p$ is the linear map
\[
dF_p:T_pM\longrightarrow T_{F(p)}N,
\qquad dF_p(v)(f)=v(f\circ F).
\]
Indeed, for $f,g\in C^\infty(N)$,
\[
dF_p(v)(fg)
=f(F(p))dF_p(v)(g)+g(F(p))dF_p(v)(f),
\]
so $dF_p(v)$ is a derivation at $F(p)$.
\end{definition}

\begin{proposition}[Properties of Differentials]
\label{prop:3-6}
\lean{diffeomorph_mfderiv_symm_eq_symm}
\leanok


\dcref{ch3:3.6}
Let $F:M\to N$ and $G:N\to P$ be smooth maps between smooth manifolds,
possibly with boundary, and let $p\in M$.
\begin{enumerate}
\item $dF_p:T_pM\to T_{F(p)}N$ is linear.
\item $d(G\circ F)_p=dG_{F(p)}\circ dF_p$.
\item $d(\operatorname{Id}_M)_p=\operatorname{Id}_{T_pM}$.
\item If $F$ is a diffeomorphism, then $dF_p$ is an isomorphism and
\[
(dF_p)^{-1}=d(F^{-1})_{F(p)}.
\]
\end{enumerate}
\end{proposition}


\begin{proposition}
\label{prop:3-8}
\lean{PointDerivation.congr_of_eqOn_nhds}
\uses{lem:3-4}
\leanok



\dcref{ch3:3.8}
Let $M$ be a smooth manifold, possibly with boundary, $p\in M$, and
$v\in T_pM$. If $f,g\in C^\infty(M)$ agree on a neighborhood of $p$, then
$v(f)=v(g)$.
\end{proposition}

\begin{proof}
\leanok
\sketch
Set $h=f-g$. Choose a smooth bump function $\psi$ supported in
$M\smallsetminus\{p\}$ and equal to $1$ on $\operatorname{supp}h$. Then
$h=\psi h$ and $h(p)=\psi(p)=0$. Lemma~\ref{lem:3-4} gives $v(h)=v(\psi h)=0$.
\end{proof}

\begin{proposition}[The Tangent Space to an Open Submanifold]
\label{prop:3-9}
\lean{mfderiv_open_subset_inclusion_isInvertible}
\uses{prop:3-8}
\leanok



\dcref{ch3:3.9}
Let $U$ be an open subset of a smooth manifold $M$, possibly with boundary,
and let $\iota:U\hookrightarrow M$ be inclusion. For every $p\in U$,
\[
d\iota_p:T_pU\longrightarrow T_pM
\]
is an isomorphism.
\end{proposition}

\begin{proof}
\leanok
\sketch
Choose a neighborhood $B$ of $p$ with $\overline B\subseteq U$. Every
$f\in C^\infty(U)$ has a smooth extension $\widetilde f\in C^\infty(M)$ that
agrees with $f$ on $\overline B$. If $d\iota_p(v)=0$, locality from
Proposition~\ref{prop:3-8} yields
\[
v(f)=v(\widetilde f|_U)=d\iota_p(v)(\widetilde f)=0,
\]
so $d\iota_p$ is injective. Conversely, for $w\in T_pM$ define
$v(f)=w(\widetilde f)$. Proposition~\ref{prop:3-8} makes this independent of the chosen
extension; it is a derivation and $d\iota_p(v)=w$.
\end{proof}

\begin{proposition}[Dimension of the Tangent Space]
\label{prop:3-10}
\lean{tangentSpace_finrank_eq_of_n_dimensional_manifold}
\uses{prop:3-6, prop:3-9}
\leanok



\dcref{ch3:3.10}
If $M$ is a smooth $n$-manifold, then $\dim T_pM=n$ for every $p\in M$.
\end{proposition}

\begin{proof}
\leanok
\sketch
For a chart $\varphi:U\to\widehat U\subseteq\mathbb{R}^n$ containing $p$,
Propositions~\ref{prop:3-6} and~\ref{prop:3-9} give
\[
T_pM\cong T_pU\cong T_{\varphi(p)}\widehat U
\cong T_{\varphi(p)}\mathbb{R}^n.
\]
The last space has dimension $n$.
\end{proof}

\begin{lemma}
\label{lem:3-11}
\lean{mfderiv_half_space_inclusion_isInvertible}
\leanok


\dcref{ch3:3.11}
For the inclusion $\iota:\mathbb{H}^n\hookrightarrow\mathbb{R}^n$ and every
$a\in\partial\mathbb{H}^n$, the map
$d\iota_a:T_a\mathbb{H}^n\to T_a\mathbb{R}^n$ is an isomorphism.
\end{lemma}

\begin{proof}
\leanok
\sketch
Every $f\in C^\infty(\mathbb{H}^n)$ admits a smooth extension
$\widetilde f$ to $\mathbb{R}^n$. Thus $d\iota_a(v)=0$ implies
$v(f)=d\iota_a(v)(\widetilde f)=0$, proving injectivity. Given
$w\in T_a\mathbb{R}^n$, set $v(f)=w(\widetilde f)$. In standard coordinates,
if $w=w^i\partial/\partial x^i|_a$, then
\[
v(f)=w^i\frac{\partial\widetilde f}{\partial x^i}(a).
\]
Continuity determines these derivatives from the restriction to
$\mathbb{H}^n$, so $v$ is well defined, is a derivation, and satisfies
$d\iota_a(v)=w$.
\end{proof}

\begin{proposition}[Dimension of Tangent Spaces on a Manifold with Boundary]
\label{prop:3-12}
\lean{tangentSpace_finrank_eq_of_n_dimensional_manifold}
\uses{lem:3-11, prop:3-6, prop:3-9}
\leanok



\dcref{ch3:3.12}
If $M$ is a smooth $n$-manifold with boundary, then $\dim T_pM=n$ for every
$p\in M$.
\end{proposition}

\begin{proof}
\leanok
\sketch
At an interior point, apply Proposition~\ref{prop:3-9} to the open submanifold
$\operatorname{Int}M$. At a boundary point, a boundary chart
$\varphi:U\to\widehat U\subseteq\mathbb{H}^n$,
Propositions~\ref{prop:3-6} and~\ref{prop:3-9}, and
Lemma~\ref{lem:3-11} give
\[
T_pM\cong T_pU\cong T_{\varphi(p)}\widehat U
\cong T_{\varphi(p)}\mathbb{H}^n
\cong T_{\varphi(p)}\mathbb{R}^n.
\]
\end{proof}

\begin{definition}
\label{def:3-14-extra-2}
\lean{lineDeriv}

\mathlibok
For a finite-dimensional vector space $V$, $a\in V$, and $v\in V$, define
\[
D_v|_a(f)=\left.\frac{d}{dt}\right|_{t=0}f(a+tv),
\qquad f\in C^\infty(V). \tag{3.5}
\]
\end{definition}

\begin{proposition}[The Tangent Space to a Vector Space]
\label{prop:3-13}
\lean{mfderiv_vector_space_to_tangent_space}
\leanok


\dcref{ch3:3.13}
For a finite-dimensional vector space $V$ with its standard smooth structure
and $a\in V$, the map $v\mapsto D_v|_a$ is a canonical isomorphism
$V\to T_aV$. It is natural with respect to linear maps: if $L:V\to W$, then
\[
dL_a(D_v|_a)=D_{Lv}|_{L(a)}. \tag{3.6}
\]
\end{proposition}

\begin{proof}
\leanok
\sketch
After choosing a basis, the Euclidean-space argument shows that $D_v|_a$ is
a derivation and that $v\mapsto D_v|_a$ is an isomorphism. For linear $L$,
\[
dL_a(D_v|_a)(f)
=\left.\frac{d}{dt}\right|_{0}f(L(a+tv))
=\left.\frac{d}{dt}\right|_{0}f(L(a)+tL(v))
=D_{Lv}|_{L(a)}(f).
\]
\end{proof}

\begin{proposition}[The Tangent Space to a Product Manifold]
\label{prop:3-14}
\lean{tangentSpaceFiniteProductEquivDirectSum_apply}
\uses{prob:3-2}
\leanok


\dcref{ch3:3.14}
Let $M_1,\ldots,M_k$ be smooth manifolds, let
$\pi_j:M_1\times\cdots\times M_k\to M_j$ be projection, and write
$p=(p_1,\ldots,p_k)$. The map
\[
\alpha:T_p(M_1\times\cdots\times M_k)
\longrightarrow T_{p_1}M_1\oplus\cdots\oplus T_{p_k}M_k,
\]
\[
\alpha(v)=(d(\pi_1)_p(v),\ldots,d(\pi_k)_p(v)) \tag{3.7}
\]
is an isomorphism. The conclusion remains valid when one factor is a smooth
manifold with boundary.
\end{proposition}

\begin{proof}
\leanok
\sketch In product charts, $\alpha$ sends a tangent vector to the blocks of its
coordinate components. Concatenating arbitrary blocks from the factor tangent
spaces defines the inverse linear map.
\end{proof}


\section{Computations in Coordinates}

\begin{definition}
\label{def:3-15-extra-1}
\lean{mfderiv_chart_coordinate_vectors_basis}
\uses{prop:3-6, prop:3-9}
\leanok



Let $(U,\varphi)$ be a smooth chart on an $n$-manifold $M$, with coordinate
functions $(x^i)$, and let $p\in U$. The coordinate vectors at $p$ are
\[
\left.\frac{\partial}{\partial x^i}\right|_p
=(d\varphi_p)^{-1}
 \left(\left.\frac{\partial}{\partial x^i}\right|_{\varphi(p)}\right)
=d(\varphi^{-1})_{\varphi(p)}
 \left(\left.\frac{\partial}{\partial x^i}\right|_{\varphi(p)}\right). \tag{3.8}
\]
For $f\in C^\infty(U)$, $\widehat f=f\circ\varphi^{-1}$ and
$\widehat p=\varphi(p)$, their action is
\[
\left.\frac{\partial}{\partial x^i}\right|_p(f)
=\frac{\partial\widehat f}{\partial x^i}(\widehat p).
\]
They form the coordinate basis associated with the chart. For a boundary
chart, use the identification
$T_{\varphi(p)}\mathbb{H}^n\cong T_{\varphi(p)}\mathbb{R}^n$ from the
inclusion differential; at a boundary point the $n$th coordinate derivative
is one-sided.
\end{definition}

\begin{proposition}
\label{prop:3-15}
\lean{chart_coordinate_vectors_basis}
\leanok


\dcref{ch3:3.15}
Let $M$ be a smooth $n$-manifold, possibly with boundary, and let $p\in M$.
For every smooth chart $(U,(x^i))$ containing $p$, the vectors
\[
\left.\frac{\partial}{\partial x^1}\right|_p,\ldots,
\left.\frac{\partial}{\partial x^n}\right|_p
\]
form a basis of the $n$-dimensional space $T_pM$.
\end{proposition}

\begin{definition}
\label{def:3-15-extra-2}
\lean{chart_coordinate_vectors_basis}
\leanok


Relative to a coordinate basis, every $v\in T_pM$ has a unique expression
\[
v=\left.v^i\frac{\partial}{\partial x^i}\right|_p.
\]
The scalars $(v^1,\ldots,v^n)$ are its components, with summation understood,
and satisfy
\[
v^j=v(x^j),
\qquad
v(x^j)=v^i\frac{\partial x^j}{\partial x^i}(p)=v^j.
\]
\end{definition}

\begin{remark}
\label{rem:3-15-extra-3}
\lean{mfderiv_eq_fderiv}

\mathlibok
Let $F:U\to V$ be smooth, where $U\subseteq\mathbb{R}^n$ and
$V\subseteq\mathbb{R}^m$ are open, with domain coordinates $(x^i)$ and
codomain coordinates $(y^j)$. The chain rule gives
\[
dF_p\left(\left.\frac{\partial}{\partial x^i}\right|_p\right)
=\left.\frac{\partial F^j}{\partial x^i}(p)
  \frac{\partial}{\partial y^j}\right|_{F(p)}. \tag{3.9}
\]
Thus the matrix of $dF_p$ in the standard coordinate bases is the Jacobian
matrix $(\partial F^j/\partial x^i)(p)$, hence agrees with the total derivative
$DF(p)$. The same holds for open subsets of half-spaces.

More generally, for charts $(U,\varphi)$ at $p$ and $(V,\psi)$ at $F(p)$,
put $\widehat F=\psi\circ F\circ\varphi^{-1}$ and
$\widehat p=\varphi(p)$. Then
\[
dF_p\left(\left.\frac{\partial}{\partial x^i}\right|_p\right)
=\left.\frac{\partial\widehat F^j}{\partial x^i}(\widehat p)
  \frac{\partial}{\partial y^j}\right|_{F(p)}. \tag{3.10}
\]
Consequently, the coordinate matrix of $dF_p$ is the Jacobian of the
coordinate representative of $F$. The notation $dF_p$ denotes this
differential, while $DF(p)$ is reserved for the total derivative between
finite-dimensional vector spaces; the differential is also called the tangent
map or pointwise pushforward.
\end{remark}

\begin{remark}
\label{rem:3-15-extra-4}
\lean{tangentCoordChange_def}

\mathlibok
Let $(U,\varphi)$ and $(V,\psi)$ be charts containing $p$, with coordinates
$(x^i)$ and $(\widetilde x^i)$, and put $\widehat p=\varphi(p)$. For the
transition map
\[
\psi\circ\varphi^{-1}(x)
=(\widetilde x^1(x),\ldots,\widetilde x^n(x)),
\]
the coordinate bases satisfy
\[
\left.\frac{\partial}{\partial x^i}\right|_p
=\left.\frac{\partial\widetilde x^j}{\partial x^i}(\widehat p)
  \frac{\partial}{\partial\widetilde x^j}\right|_p. \tag{3.11}
\]
Hence, if
$v=v^i\partial/\partial x^i|_p
=\widetilde v^j\partial/\partial\widetilde x^j|_p$, then
\[
\widetilde v^j
=\frac{\partial\widetilde x^j}{\partial x^i}(\widehat p)v^i. \tag{3.12}
\]
\end{remark}

\begin{example}
\label{ex:3-16}
\lean{polar_vector_in_standard_coordinates}
\leanok


\dcref{ch3:3.16}
For $(x,y)=(r\cos\theta,r\sin\theta)$, let $p$ have polar coordinates
$(2,\pi/2)$. At $p$,
\[
\left.\frac{\partial}{\partial r}\right|_p
=\left.\frac{\partial}{\partial y}\right|_p,
\qquad
\left.\frac{\partial}{\partial\theta}\right|_p
=-2\left.\frac{\partial}{\partial x}\right|_p.
\]
Therefore the vector
\[
v=3\left.\frac{\partial}{\partial r}\right|_p
-\left.\frac{\partial}{\partial\theta}\right|_p
\]
has standard-coordinate expression
\[
v=2\left.\frac{\partial}{\partial x}\right|_p
+3\left.\frac{\partial}{\partial y}\right|_p.
\]
\end{example}

\begin{remark}
\label{rem:3-15-extra-5}
\lean{tangentMap_chart_symm}

\mathlibok
The vector $\partial/\partial x^i|_p$ depends on the entire coordinate system,
because its definition holds all other coordinates constant; changing those
coordinates can change the vector even when the function $x^i$ itself is
unchanged.
\end{remark}



\section{The Tangent Bundle}

\begin{definition}
\label{def:3-16-extra-1}
\lean{TangentBundle}

\mathlibok
For a smooth manifold $M$, possibly with boundary, its tangent bundle is
\[
TM=\coprod_{p\in M}T_pM.
\]
An element may be written $(p,v)$, $v_p$, or $v\in T_pM$. The bundle
projection is
\[
\pi:TM\to M,qquad \pi(p,v)=p,
\]
and each $T_pM$ is identified with its canonical image in $TM$.
\end{definition}

\begin{proposition}
\label{prop:3-18}
\lean{Bundle.contMDiff_proj}
\uses{lem:1-35}

\mathlibok

\dcref{ch3:3.18}
For a smooth $n$-manifold $M$, the set $TM$ has a natural topology and smooth
structure making it a smooth $2n$-manifold, and $\pi:TM\to M$ is smooth.
\end{proposition}

\begin{proof}
\sketch
For a chart $(U,\varphi)$ with coordinate functions $(x^i)$, define
\[
\widetilde\varphi:\pi^{-1}(U)\to\varphi(U)\times\mathbb{R}^n,
\qquad
\widetilde\varphi
\left(\left.v^i\frac{\partial}{\partial x^i}\right|_p\right)
=(x^1(p),\ldots,x^n(p),v^1,\ldots,v^n). \tag{3.13}
\]
This is a bijection with inverse
\[
(x^1,\ldots,x^n,v^1,\ldots,v^n)
\longmapsto
\left.v^i\frac{\partial}{\partial x^i}\right|_{\varphi^{-1}(x)}.
\]
On an overlap with coordinates $(\widetilde x^i)$, (3.12) gives the transition
map
\[
(x,v)\longmapsto
\left(\widetilde x(x),
\frac{\partial\widetilde x^1}{\partial x^j}(x)v^j,\ldots,
\frac{\partial\widetilde x^n}{\partial x^j}(x)v^j\right),
\]
which is smooth. A countable coordinate cover and the smooth manifold chart
lemma produce the topology and smooth structure. Points in a common fiber lie
in a common bundle chart; points over distinct base points are separated by
inverse images of disjoint base charts, proving Hausdorffness. In these charts,
$\pi(x,v)=x$, so $\pi$ is smooth.
\end{proof}

\begin{notation}
\label{not:3-16-extra-2}
\lean{chartAt}

\mathlibok
The coordinates $(x^i,v^i)$ in (3.13) are the natural coordinates on $TM$.
\end{notation}

\begin{proposition}[Tangent Bundles of Manifolds with Boundary]
\label{exer:3-19}
\lean{tangentBundleBoundaryChart_apply}
\leanok


\dcref{ch3:3.19}
If $M$ is a smooth manifold with boundary, its natural topology and smooth
structure make $TM$ a smooth manifold with boundary. A boundary chart
$(U,(x^i))$ on $M$ induces the boundary chart
$(\pi^{-1}(U),(v^i,x^i))$ by reordering the natural coordinates
$(x^i,v^i)$.
\end{proposition}

\begin{definition}
\label{def:3-16-extra-3}
\lean{tangentMap}

\mathlibok
For a smooth map $F:M\to N$, the global differential, or global tangent map,
is
\[
dF:TM\to TN,
\qquad dF|_{T_pM}=dF_p.
\]
Thus $dF(v)=dF_p(v)$ for $v\in T_pM$.
\end{definition}

\begin{proposition}
\label{prop:3-21}
\lean{fun}

\mathlibok
\dcref{ch3:3.21}
If $F:M\to N$ is smooth, then $dF:TM\to TN$ is smooth.
\end{proposition}

\begin{proof}
\sketch
In natural coordinates, (3.9) gives
\[
dF(x,v)=
\left(F(x),
\left(\frac{\partial F^\alpha}{\partial x^i}(x)v^i\right)_\alpha\right).
\]
This map is smooth because the coordinate functions of $F$ and their first
partial derivatives are smooth.
\end{proof}

\begin{notation}
\label{not:3-16-extra-4}
\lean{diffeomorph_mfderiv_symm_eq_symm}
\leanok


For a diffeomorphism $F$, write $dF^{-1}$ for either $(dF)^{-1}$ or
$d(F^{-1})$; these maps agree.
\end{notation}


\section{Velocity Vectors of Curves}

\begin{definition}
\label{def:3-17-extra-1}
\lean{curve_velocity}
\leanok


A curve in a manifold $M$, possibly with boundary, is a continuous map
$\gamma:J\to M$ from an interval $J\subseteq\mathbb{R}$. Curves are
parametrized. If $\gamma$ is smooth and $t_0\in J$, its velocity is
\[
\gamma'(t_0)
=d\gamma\left(\left.\frac{d}{dt}\right|_{t_0}\right)
\in T_{\gamma(t_0)}M.
\]
It is also denoted $\dot\gamma(t_0)$ or $d\gamma/dt|_{t=t_0}$ and acts by
\[
\gamma'(t_0)(f)=\left.\frac{d}{dt}\right|_{t_0}(f\circ\gamma).
\]
At an endpoint of $J$, derivatives are one-sided, equivalently computed from
a smooth local extension. In coordinates $(x^i)$, writing the coordinate
functions of the curve as $\gamma^i$, one has
\[
\gamma'(t_0)
=\left.\frac{d\gamma^i}{dt}(t_0)
  \frac{\partial}{\partial x^i}\right|_{\gamma(t_0)}.
\]
\end{definition}

\begin{proposition}
\label{prop:3-23}
\lean{exists_smooth_curve_with_velocity_boundaryless}
\leanok


\dcref{ch3:3.23}
Let $M$ be a smooth manifold, possibly with boundary, and let $p\in M$. Every
$v\in T_pM$ is the velocity of a smooth curve in $M$.
\end{proposition}

\begin{proof}
\leanok
\sketch
Choose coordinates centered at $p$ and write
$v=v^i\partial/\partial x^i|_p$. For sufficiently small $\varepsilon>0$, use
the coordinate curve
\[
\gamma(t)=(tv^1,\ldots,tv^n). \tag{3.14}
\]
At an interior point its domain is $(-\varepsilon,\varepsilon)$. At a boundary
point, take this domain when $v^n=0$, $[0,\varepsilon)$ when $v^n>0$, and
$(-\varepsilon,0]$ when $v^n<0$. In every case $\gamma(0)=p$ and
$\gamma'(0)=v$.
\end{proof}

\begin{proposition}[The Velocity of a Composite Curve]
\label{prop:3-24}
\lean{composite_curve_velocity_of_contMDiff}
\leanok


\dcref{ch3:3.24}
Let $F:M\to N$ be smooth and $\gamma:J\to M$ a smooth curve. For every
$t_0\in J$,
\[
(F\circ\gamma)'(t_0)=dF(\gamma'(t_0)).
\]
\end{proposition}

\begin{proof}
\leanok
\sketch
By the chain rule for differentials,
\[
(F\circ\gamma)'(t_0)
=d(F\circ\gamma)\left(\left.\frac{d}{dt}\right|_{t_0}\right)
=dF\left(d\gamma\left(\left.\frac{d}{dt}\right|_{t_0}\right)\right).
\]
\end{proof}

\begin{corollary}[Computing the Differential Using a Velocity Vector]
\label{cor:3-25}
\lean{differential_eq_velocity_of_composite_curve}
\uses{prop:3-23, prop:3-24}
\leanok



\dcref{ch3:3.25}
Let $F:M\to N$ be smooth, $p\in M$, and $v\in T_pM$. For every smooth curve
$\gamma:J\to M$ with $0\in J$, $\gamma(0)=p$, and $\gamma'(0)=v$,
\[
dF_p(v)=(F\circ\gamma)'(0).
\]
\end{corollary}


\section{Alternative Definitions of the Tangent Space}

\begin{definition}
\label{def:3-18-extra-1}
\lean{smoothSheafCommRing}

\mathlibok
A smooth function element on $M$ is a pair $(f,U)$ with $U\subseteq M$ open
and $f:U\to\mathbb{R}$ smooth. For $p\in M$, restrict to function elements
whose domains contain $p$, and declare
\[
(f,U)\sim(g,V)
\quad\Longleftrightarrow\quad
f|_W=g|_W\text{ for some open neighborhood }W\subseteq U\cap V\text{ of }p.
\]
The equivalence class $[f]_p$ is the germ of $f$ at $p$, and the algebra of
all such germs is denoted $C_p^\infty(M)$. Its operations are
\[
c[f]_p=[cf]_p,qquad [f]_p+[g]_p=[f+g]_p,
\qquad [f]_p[g]_p=[fg]_p,
\]
where the last two representatives are restricted to $U\cap V$. The zero
element is the germ of the zero function.
\end{definition}

\begin{definition}
\label{def:3-18-extra-4}
\lean{IsCoordinateTangentVector}
\leanok


Equivalently, a tangent vector at $p$ can be represented by an assignment of a
tuple $(v^1,\ldots,v^n)\in\mathbb{R}^n$ to every smooth coordinate chart
containing $p$, such that on overlaps the tuples obey the transformation law
(3.12).
\end{definition}


\section{Categories and Functors}

\begin{definition}
\label{def:3-19-extra-1}
\lean{Category}

\mathlibok
A category $\mathrm C$ consists of classes $\operatorname{Ob}(\mathrm C)$ and
$\operatorname{Hom}(\mathrm C)$, source and target objects for each morphism,
and composition maps
\[
\operatorname{Hom}_{\mathrm C}(X,Y)\times
\operatorname{Hom}_{\mathrm C}(Y,Z)
\longrightarrow\operatorname{Hom}_{\mathrm C}(X,Z),
\qquad (f,g)\longmapsto g\circ f.
\]
Composition is associative, and every object $X$ has an identity morphism
$\operatorname{Id}_X$ satisfying
\[
\operatorname{Id}_Y\circ f=f=f\circ\operatorname{Id}_X
\qquad(f\in\operatorname{Hom}_{\mathrm C}(X,Y)).
\]
\end{definition}

\begin{definition}
\label{def:3-19-extra-2}
\lean{CategoryTheory.IsIso}

\mathlibok
A morphism $f\in\operatorname{Hom}_{\mathrm C}(X,Y)$ is an isomorphism if
there is $g\in\operatorname{Hom}_{\mathrm C}(Y,X)$ such that
\[
f\circ g=\operatorname{Id}_Y,
\qquad g\circ f=\operatorname{Id}_X.
\]
\end{definition}

\begin{example}[Categories]
\label{ex:3-26}
\lean{pointed_map_iff_exists_hom}
\leanok


\dcref{ch3:3.26}
Examples of categories include sets and maps; topological spaces or
manifolds and continuous maps; smooth manifolds, with or without boundary, and
smooth maps; real or complex vector spaces and linear maps; groups or abelian
groups and homomorphisms; and rings or commutative rings and homomorphisms.

A pointed object is a pair $(X,p)$ consisting of an object and a distinguished
point. A morphism
\[
F:(X,p)\longrightarrow(X',p')
\]
is pointed when $F(p)=p'$. Pointed sets, topological spaces, topological
manifolds, and smooth manifolds with their pointed maps form corresponding
categories.
\end{example}

\begin{definition}
\label{def:3-19-extra-3}
\lean{LocallySmall}

\mathlibok
A category is small if both its object class and its morphism class are sets.
It is locally small if $\operatorname{Hom}_{\mathrm C}(X,Y)$ is a set for every
pair of objects $X,Y$. Each category in Example~\ref{ex:3-26} is locally small
but not small.
\end{definition}

\begin{definition}
\label{def:3-19-extra-4}
\lean{CategoryTheory.Functor}

\mathlibok
For categories $\mathrm C$ and $\mathrm D$, a covariant functor
$\mathcal F:\mathrm C\to\mathrm D$ assigns an object $\mathcal F(X)$ to every
object $X$ and a morphism
\[
\mathcal F(f)\in
\operatorname{Hom}_{\mathrm D}(\mathcal F(X),\mathcal F(Y))
\]
to every $f\in\operatorname{Hom}_{\mathrm C}(X,Y)$, preserving identities and
composition:
\[
\mathcal F(\operatorname{Id}_X)=\operatorname{Id}_{\mathcal F(X)},
\qquad
\mathcal F(g\circ h)=\mathcal F(g)\circ\mathcal F(h).
\]
\end{definition}

\begin{exercise}
\label{exer:3-17}
\lean{cubicShear_symm_fderiv_apply_e1}
\leanok
\dcref{ch3:3.17}
Let $\left( {x, y}\right)$ denote the standard coordinates on ${\mathbb{R}}^{2}$ . Verify that $\left( {\widetilde{x},\widetilde{y}}\right)$ are global smooth coordinates on ${\mathbb{R}}^{2}$ , where

\[
\widetilde{x} = x,\;\widetilde{y} = y + {x}^{3}.
\]

Let $p$ be the point $\left( {1,0}\right)  \in  {\mathbb{R}}^{2}$ (in standard coordinates), and show that

\[
{\left. \frac{\partial }{\partial x}\right| }_{p} \neq  {\left. \frac{\partial }{\partial \widetilde{x}}\right| }_{p}
\]

even though the coordinate functions $x$ and $\widetilde{x}$ are identically equal.
\end{exercise}

\begin{exercise}
\label{exer:3-27}
\lean{Functor.map_isIso}
\mathlibok
\dcref{ch3:3.27}
Show that any (covariant or contravariant) functor from C to D takes isomorphisms in C to isomorphisms in D.
\end{exercise}

\begin{exercise}
\label{exer:3-5}
\lean{tangent_vector_mul_eq_zero_of_vanish_at}
\leanok
\uses{lem:3-4}
\dcref{ch3:3.5}
\begin{itemize}
\item Exercise 3.5. Prove Lemma 3.4.
\end{itemize}
\end{exercise}

\begin{exercise}
\label{exer:3-7}
\lean{mfderiv}
\mathlibok
\uses{prop:3-6}
\dcref{ch3:3.7}
Prove Proposition 3.6.
\end{exercise}
